% !TEX root = ../ipeesec.tex
\section{Modelling of IPsec}\label{sec:modelling}

\subsection{Key exchange}\label{sec:key-exchange}

Regarding the key exchange using elliptic curve cryptography, the reader is invited to read through~\cite{happelEllipticCurvesF42022} for a quick introduction.
Crucially, we will assume that a shared encryption algorithm $E_{p,q} : T \to S$ for a private key $p$ and a public key $q$ is chosen, together with a decryption algorithm $D_{p,q} : S \to T$ such that $D_{p,q} \circ E_{p,q} = \id_T$, making $D_{p,q}$ a retraction of the section $E_{p,q}$.

After successfully establishing public and private keys and an encryption/decryption algorithm, we will model this as two tuples $(p_c, q_c), (p_s, q_s) : K^2$ for the private and public keys of the client and of the server respectively, together with $E : \Hom(K^2, \hom(T, S))$ and $D : \Hom(K^2, \hom(S, T))$.

\subsection{Authentication}\label{sec:authentication}

\todo[inline]{Needs to be written.}

\subsection{The tunnel}\label{sec:tunnel}

We are left modelling the traffic selector, which we will model as two relations $Q_c, Q_s \rightarrowtail \mathrm{IP\text{-}Range}$ for the client and server respectively, where
\[
	\mathrm{IP\text{-}Range} \cong [0, 2^{32} - 1] \cong [0, 2^{8} - 1]^4,
	\qquad
	a_1 \cdot 2^{24} + a_2 \cdot 2^{16} + a_3 \cdot 2^{8} + a_4 \;\mapsfrom\; (a_1, a_2, a_3, a_4),
\]
is commonly represented as
\[
	\mathrm{repr}_{\mathrm{IP}} : \mathrm{IP\text{-}Range} \to \mathrm{String},
	\qquad
	\mathrm{repr}_{\mathrm{IP}}(a_1, a_2, a_3, a_4) = a_1 \texttt{.} a_2 \texttt{.} a_3 \texttt{.} a_4,
\]
with each $a_i$ written in decimal, for example $\mathrm{repr}_{\mathrm{IP}}(10, 0, 1, 0) = \texttt{10.0.1.0}$.
We can trivially convert these subobjects to their respective characteristic functions through the topos structure on $\Set$, yielding
\[
	\chi_{Q_c}, \chi_{Q_s} : \mathrm{IP\text{-}Range} \to \{\top, \bot\} \cong \Bool
\]
for the client and server respectively, which are exactly the classical filters.

We will model the communication layer by first considering only streaming data from one party to the other without responses.
For this, the topos of trees~\cite{birkedalFirstStepsSynthetic2012} is a natural semantic choice that provides us with a rich, well-known syntax.
Crucially, by use of the later modality $\later$, we can encode the passage of time on the typing level, where each discrete time step in the communication is given by an ordinal number $t \in \omega = \{0 \le 1 \le 2 \le \dots\}$.

Consider $\Str(A) = A \times \later \Str(A)$ for some data $A$ that the client intends to transmit.
From now on, we will only consider the case of the client communicating to the server, as the other way around follows trivially.
We first need a productive stream transformer, which obviously is given by lifting a morphism $f : A \to B$, where $B = \im E_{p_c, q_s}(A)$, first through the functoriality of the diagonal functor $\Delta : \Set \to \toposoftrees$ and then through the functoriality of $\Str$, yielding
\[
	\Str(\Delta f) : \Str(\Delta A) \to \Str(\Delta B).
\]
We next need a shim
\[
	\mathrm{shim} : \Str(B) \to \Str(B + 1),
\]
which represents the UDP transport layer, encoding possible data loss by use of the optional data type $B + 1$.
Finally, the recipient can then decode the message (lossily) by lifting $D_{p_s, q_c}$ analogously to
\[
	\Str(D_{p_s, q_c} + 1) : \Str(B + 1) \to \Str(A + 1).
\]

How can we eliminate this loss, one may ask?
This paves the way for us to consider \emph{bidirectional communication}.
We will model this by passing over from $\omega$ to the first uncountable ordinal $\omega_1$, corresponding to the topos $\Set^{\omega_1^{\mathrm{op}}}$ of presheaves on $\omega_1$, and introducing a map
\[
	\mathrm{skip} : \omega_1 \to \omega_1, \qquad \mathrm{skip}(n) = \omega + n.
\]
This then allows us to encode finishing one stream transfer and having the recipient respond.
We may in fact respond eagerly by performing $\mathrm{skip}$ at any given $t \in \omega_1$ which is not a limit ordinal.
